\documentclass[11pt,a4paper]{article}
\usepackage[margin=1in]{geometry}
\usepackage{hyperref}
\usepackage{enumitem}
\usepackage{graphicx}
\usepackage{xcolor}

% -----------------------------------------------------
% bvraghav-tiet-ucs503p-article.sty
% -----------------------------------------------------

% Variable: Lab Instructor
\makeatletter \newcommand{\BvrSetLabInstructor}[1]{%
  \gdef\Bvr@LabInstructor{#1}}
% default
\newcommand{\Bvr@LabInstructor}{Dr. Raghav %
  B. Venkataramaiyer}
% public accessor
\newcommand{\BvrLabInstructorValue}{\Bvr@LabInstructor}
\makeatother

% Add subtitle
\usepackage{titling}
% -- Set title bold
\pretitle{\begin{center}\bfseries\fontsize{18bp}{18bp}\selectfont}
\makeatletter
\newcommand{\BvrSubtitle}[1]{%
  \posttitle{%
    \par\end{center}
    \begin{center}\large#1\end{center}
    \vskip 0.5em}%
}
\makeatother

% Hints in the section title(s)
\newcommand{\BvrHint}[1]{%
  {\normalfont\normalsize\itshape\hfill #1}}

% -----------------------------------------------------

\title{AlgoRecall}

\BvrSetLabInstructor{Ms. Paramveer Kaur}

\BvrSubtitle{%
  Adaptive DSA Revision and Analytics Platform \\
  Project Proposal for UCS503P  \\
  Submitted to: \BvrLabInstructorValue \\ \bfseries
  Thapar Institute of Engineering and Technology%
}

\author{
  Pihu Tanoch, CSED%
  \thanks{Roll No: \texttt{1024170373}}
  \and
  Ekmanjot Kaur, CSED%
  \thanks{Roll No: \texttt{1024170377}}
  \and
  Mannatpreet Singh, CSED%
  \thanks{Roll No: \texttt{1024170382}}
}

\date{\today}

\begin{document}
\maketitle

\section{Abstract}

AlgoRecall will be a web-based platform designed to help users manage and
organize the revision of previously solved DSA questions. The system will
allow users to store solved questions along with details such as topic,
difficulty level, confidence level, and revision status.

The platform will maintain a complete revision history and schedule future
revisions based on user-defined and system-generated revision parameters.
It will provide a dedicated dashboard for viewing due questions and
tracking revision activities. Users will be able to search, filter, and
categorize questions to efficiently manage their revision workflow.

In addition, the system will offer analytical reports and visualizations
that will summarize revision progress, topic-wise performance, confidence
distribution, and revision consistency.

\section{Introduction}

Data Structures and Algorithms (DSA) play a crucial role in technical
interviews and competitive programming. As part of their preparation,
students and job aspirants solve a large number of problems across various
topics such as Arrays, Linked Lists, Trees, Graphs, Dynamic Programming,
and Greedy Algorithms. Continuous practice helps learners develop
problem-solving skills and understand different algorithmic patterns.

As the number of solved questions grows, managing and revisiting
previously learned concepts becomes increasingly important. Learners often
maintain records of solved problems using spreadsheets, notes, or online
coding platforms. However, these methods primarily focus on problem-solving
activities and provide limited support for organizing and tracking
revisions.

To address this need, AlgoRecall is proposed as a web-based platform that
will enable users to maintain a structured record of solved questions,
track revision activities, and monitor their overall learning progress
through various management and analytical features.

\section{Problem Statement}

Existing coding platforms and online judges primarily focus on helping
users solve programming problems and evaluate their coding skills. While
these platforms provide features such as problem categorization, progress
tracking, and submission history, they offer limited support for systematic
revision and long-term knowledge retention.

As learners solve an increasing number of DSA questions, it becomes
difficult to remember previously learned concepts, problem-solving
patterns, and algorithmic approaches. Revisions are often performed
irregularly and without a structured plan, making it challenging for
learners to determine which questions require more attention.
Additionally, identifying weak topics and monitoring revision progress
becomes increasingly difficult as the volume of solved questions grows.

Therefore, there is a need for a dedicated revision management system that
can organize solved questions, schedule revisions effectively, maintain
revision records, and provide meaningful insights into a learner's
strengths and weaknesses. Such a system can support a more disciplined and
structured approach to DSA revision and learning.

\section{Objectives}

The primary objective of AlgoRecall is to provide a structured platform
for managing and revising previously solved DSA questions. The system will
aim to support users in maintaining consistent revision habits and
tracking their learning progress effectively.

The objectives of the project are:

\begin{itemize}[leftmargin=0.7cm]
    \item To store and manage records of solved DSA questions.
    \item To allow users to assign and update confidence levels for each question.
    \item To schedule and organize future revisions based on revision-related parameters.
    \item To maintain a complete history of revision activities.
    \item To identify weak topics and areas that require additional practice.
    \item To support a structured and disciplined revision workflow.
    \item To improve long-term retention of DSA concepts and problem-solving patterns.
\end{itemize}

\section{Scope of the Project}

\subsection{In Scope}

\begin{enumerate}[leftmargin=0.7cm]
    \item User Registration and Authentication
    \item Question Management (Add, Edit, Delete, View)
    \item Confidence Level Tracking
    \item Revision Scheduling
    \item Revision History
    \item Analytics and Reports
    \item Search and Filter Functionality
\end{enumerate}

\subsection{Out of Scope}

\begin{enumerate}[leftmargin=0.7cm]
    \item Online Judge Functionality
    \item Code Compilation and Execution
    \item Contest Management
    \item AI-based Solution Generation
    \item Automatic LeetCode Synchronization
    \item Real-time Collaboration Features
\end{enumerate}

The project will be focused on revision management and learning analytics
rather than problem solving or code execution.

\section{Target Users}

The primary stakeholders of AlgoRecall shall be DSA learners, students
preparing for technical interviews, and competitive programmers.

\section{Functional Requirements}

The system will provide the following functionalities:

\begin{itemize}[leftmargin=0.7cm]
    \item User Registration and Login
    \item Add, Edit, and Delete Questions
    \item Revise Questions and Record Revision Activity
    \item Update Confidence Levels
    \item View Due Questions for Revision
    \item View Analytics and Progress Reports
    \item Search and Filter Questions
\end{itemize}

\section{Non-Functional Requirements}

The system will satisfy the following quality requirements:

\begin{itemize}[leftmargin=0.7cm]
    \item \textbf{Performance:} Quick response time for user operations.
    \item \textbf{Security:} Secure authentication and protection of user data.
    \item \textbf{Usability:} Simple and user-friendly interface.
    \item \textbf{Scalability:} Ability to handle increasing amounts of data and users.
    \item \textbf{Maintainability:} Easy to update and extend in the future.
    \item \textbf{Availability:} Reliable access to the system whenever required.
\end{itemize}

\section{System Features}

The system will consist of the following major modules:

\begin{itemize}[leftmargin=0.7cm]
    \item \textbf{Authentication:} User registration, login, and secure access management.
    \item \textbf{Question Management:} Add, edit, delete, and organize solved DSA questions.
    \item \textbf{Revision Engine:} Schedule and manage future revisions.
    \item \textbf{Due Today Dashboard:} Display questions that are due for revision.
    \item \textbf{Revision History:} Maintain records of past revision activities.
    \item \textbf{Analytics:} Provide insights into revision progress and performance.
    \item \textbf{Search \& Filter:} Enable quick retrieval of questions based on different criteria.
\end{itemize}

\section{Database Design}

The system will use a MongoDB database consisting of three primary
collections:

\begin{itemize}[leftmargin=0.7cm]
    \item \textbf{Users:} User information such as name, email, password hash,
    and account-related details required for authentication and authorization.

    \item \textbf{UserQuestions:} All questions added by a user along with
    attributes such as title, platform, topic, difficulty level, confidence
    level, date solved, last revision date, next revision date, revision
    count, and current revision status.

    \item \textbf{RevisionHistory:} Records of individual revision sessions,
    including the associated question, revision date, confidence updates,
    and other revision-related information.
\end{itemize}

\section{Revision Engine}

The Revision Engine will be responsible for scheduling and managing
question revisions. It will utilize factors such as Confidence Level,
Question Difficulty, and Last Revision Date to determine when a question
should be revised next.

Based on these parameters, the system shall calculate the next revision
date and assign a priority to each question. Questions with lower
confidence levels, higher difficulty, or longer gaps since the last
revision shall be given higher priority.

\section{User Workflow}

The typical workflow of the system will be as follows:

\begin{enumerate}[leftmargin=0.7cm]
    \item The user registers or logs into the platform.
    \item The user adds solved DSA questions with details such as topic,
    difficulty, and confidence level.
    \item The system schedules future revisions for the added questions.
    \item The user views questions that are due for revision through the dashboard.
    \item The user revises a question and updates their confidence level.
    \item The system records the revision activity and reschedules the next revision.
    \item The user can monitor progress and performance through analytics and reports.
\end{enumerate}

\section{Technical Stack to be Used}

\subsection{Frontend Layer}

\begin{itemize}[leftmargin=0.7cm]
    \item React.js
    \item Tailwind CSS
    \item JavaScript (ES6+)
\end{itemize}

\subsection{Backend Layer}

\begin{itemize}[leftmargin=0.7cm]
    \item Node.js
    \item Express.js
    \item RESTful APIs
\end{itemize}

\subsection{Database Layer}

\begin{itemize}[leftmargin=0.7cm]
    \item MongoDB Atlas
    \item Mongoose ODM
\end{itemize}

\subsection{Authentication \& Security}

\begin{itemize}[leftmargin=0.7cm]
    \item JSON Web Tokens (JWT)
    \item Password Hashing (bcrypt)
\end{itemize}

\subsection{Data Visualization}

\begin{itemize}[leftmargin=0.7cm]
    \item Chart.js
\end{itemize}

\subsection{Version Control \& Development Tools}

\begin{itemize}[leftmargin=0.7cm]
    \item Git
    \item GitHub
    \item Visual Studio Code
\end{itemize}

\section{Expected Outcomes}

The project is expected to improve revision consistency, retention of DSA
concepts, identification of weak topics, and overall interview preparation.

\end{document}
